\section{Per-task results}
\label{app:pertask}
All 29 tasks over the six repetitions of the headline basis, ordered by flip
rate. \textbf{The \emph{group} column describes the published runs' behaviour,
not ours}: it records whether the benchmark's own published trajectories all
passed a task, all failed it, or disagreed on it (Section~\ref{sec:setup}), and
it is how the task set was stratified. A task can therefore be listed as
\emph{all-pass} and still have failed in every one of our repetitions.\footnote{%
Four tasks here score 0/6 despite having passed in two of the three rounds that
preceded the shift in behaviour of Section~\ref{sec:endpoint}:
\texttt{BidFileRenameTask}, \texttt{CheckEventTimeTask},
\texttt{DownloadSendReceiptTask} and \texttt{MattermostCreateChannelTask} (see
Figure~\ref{fig:dumbbell}). The headline basis lies entirely after that change,
so these rows are a consequence of it rather than an anomaly. The remaining
zero-scoring tasks are ones the published runs also fail.} Flip rate is the share of trial pairs that disagree,
$k(n-k)/\binom{n}{2}$; its maximum is below 1 and depends on $n$, reaching
0.600 at $n=6$. Wilson intervals are wide at this sample size, which is why the
suite-level figures rather than the per-task ones carry the argument in
Section~\ref{sec:noise}.

\begin{center}\scriptsize
\begin{tabular}{llccccc}
\toprule
task & group & passed & pass rate & Wilson 95\% & flip rate & steps (median) \\
\midrule
\texttt{CVEmailTask} & disagreed & 3/6 & 50\% & 19--81\% & 0.600 & 11 \\
\texttt{ScheduleLunchViaSmsTask} & disagreed & 4/6 & 67\% & 30--90\% & 0.533 & 25 \\
\texttt{MattermostBudgetApprovalPipelineTask} & disagreed & 2/6 & 33\% & 10--70\% & 0.533 & 50 \\
\texttt{MattermostEmailTask} & disagreed & 2/6 & 33\% & 10--70\% & 0.533 & 21 \\
\texttt{TakeSelfieTask} & all-pass & 5/6 & 83\% & 44--97\% & 0.333 & 12 \\
\texttt{MastodonImportMutedUsersTask} & disagreed & 1/6 & 17\% & 3--56\% & 0.333 & 50 \\
\texttt{SendFormsTask} & disagreed & 1/6 & 17\% & 3--56\% & 0.333 & 50 \\
\texttt{AcceptMeetingTask} & all-pass & 6/6 & 100\% & 61--100\% & 0.000 & 7 \\
\texttt{ChangeWallpaperTask} & all-pass & 6/6 & 100\% & 61--100\% & 0.000 & 9 \\
\texttt{CloseFlightModeTask} & all-pass & 6/6 & 100\% & 61--100\% & 0.000 & 7 \\
\texttt{GoogleMapsAlibabaSouthNeighborTask} & all-pass & 6/6 & 100\% & 61--100\% & 0.000 & 9 \\
\texttt{MastodonChangeHeaderTask} & all-pass & 6/6 & 100\% & 61--100\% & 0.000 & 7 \\
\texttt{MastodonNewPostTask} & all-pass & 6/6 & 100\% & 61--100\% & 0.000 & 6 \\
\texttt{OpenFlightModeTask} & all-pass & 6/6 & 100\% & 61--100\% & 0.000 & 6 \\
\texttt{BidFileRenameTask} & disagreed & 0/6 & 0\% & 0--39\% & 0.000 & 50 \\
\texttt{CheckEventTimeTask} & all-pass & 0/6 & 0\% & 0--39\% & 0.000 & 20 \\
\texttt{CheckInterviewTimesTask} & all-fail & 0/6 & 0\% & 0--39\% & 0.000 & 50 \\
\texttt{DownloadSendReceiptTask} & disagreed & 0/6 & 0\% & 0--39\% & 0.000 & 34 \\
\texttt{GraduationMassEmailTask} & disagreed & 0/6 & 0\% & 0--39\% & 0.000 & 13 \\
\texttt{LocalFileManagementTask2} & all-fail & 0/6 & 0\% & 0--39\% & 0.000 & 50 \\
\texttt{MastodonCalendarMultiMemosTask} & all-fail & 0/6 & 0\% & 0--39\% & 0.000 & 50 \\
\texttt{MastodonGetServerInfoTask} & all-fail & 0/6 & 0\% & 0--39\% & 0.000 & 50 \\
\texttt{MastodonServerInfoReportTask} & all-fail & 0/6 & 0\% & 0--39\% & 0.000 & 50 \\
\texttt{MastodonShareLocationTask} & all-fail & 0/6 & 0\% & 0--39\% & 0.000 & 36 \\
\texttt{MattermostCreateChannelTask} & disagreed & 0/6 & 0\% & 0--39\% & 0.000 & 24 \\
\texttt{MattermostProjectStatusReportTask} & all-fail & 0/6 & 0\% & 0--39\% & 0.000 & 50 \\
\texttt{MattermostResourceConflictResolutionTask} & all-fail & 0/6 & 0\% & 0--39\% & 0.000 & 50 \\
\texttt{ReadQwen3PaperTask5} & all-fail & 0/6 & 0\% & 0--39\% & 0.000 & 50 \\
\texttt{SMSManagement} & all-fail & 0/6 & 0\% & 0--39\% & 0.000 & 50 \\
\bottomrule
\end{tabular}

\end{center}

\section{Tasks that never flipped}
\label{app:stable}
The 22 tasks that returned the same verdict in every repetition of the headline
basis. As in Appendix~\ref{app:pertask}, \emph{group} is the published runs'
behaviour rather than ours, which is why a row can read \emph{all-pass} and
\emph{always failed} at once. Their existence is the control: a setup that flipped everywhere would be
measuring something other than the benchmark. Note that all 10 tasks the
published runs uniformly fail are here, and 8 of the 9 they uniformly pass.

\begin{center}\scriptsize
\begin{tabular}{llcl}
\toprule
task & group & passed & outcome \\
\midrule
\texttt{CheckInterviewTimesTask} & all-fail & 0/6 & always failed \\
\texttt{LocalFileManagementTask2} & all-fail & 0/6 & always failed \\
\texttt{MastodonCalendarMultiMemosTask} & all-fail & 0/6 & always failed \\
\texttt{MastodonGetServerInfoTask} & all-fail & 0/6 & always failed \\
\texttt{MastodonServerInfoReportTask} & all-fail & 0/6 & always failed \\
\texttt{MastodonShareLocationTask} & all-fail & 0/6 & always failed \\
\texttt{MattermostProjectStatusReportTask} & all-fail & 0/6 & always failed \\
\texttt{MattermostResourceConflictResolutionTask} & all-fail & 0/6 & always failed \\
\texttt{ReadQwen3PaperTask5} & all-fail & 0/6 & always failed \\
\texttt{SMSManagement} & all-fail & 0/6 & always failed \\
\texttt{AcceptMeetingTask} & all-pass & 6/6 & always passed \\
\texttt{ChangeWallpaperTask} & all-pass & 6/6 & always passed \\
\texttt{CheckEventTimeTask} & all-pass & 0/6 & always failed \\
\texttt{CloseFlightModeTask} & all-pass & 6/6 & always passed \\
\texttt{GoogleMapsAlibabaSouthNeighborTask} & all-pass & 6/6 & always passed \\
\texttt{MastodonChangeHeaderTask} & all-pass & 6/6 & always passed \\
\texttt{MastodonNewPostTask} & all-pass & 6/6 & always passed \\
\texttt{OpenFlightModeTask} & all-pass & 6/6 & always passed \\
\texttt{BidFileRenameTask} & disagreed & 0/6 & always failed \\
\texttt{DownloadSendReceiptTask} & disagreed & 0/6 & always failed \\
\texttt{GraduationMassEmailTask} & disagreed & 0/6 & always failed \\
\texttt{MattermostCreateChannelTask} & disagreed & 0/6 & always failed \\
\bottomrule
\end{tabular}

\end{center}

\section{Three intervals that are easy to confuse}
\label{app:intervals}
We compute and report three different quantities, and conflating them is the
most likely way to misread this paper.
\begin{itemize}\itemsep2pt
\item \textbf{Re-run distribution.} Where the total would land if the same task
  set were run again. Each task's rate is drawn from a Jeffreys posterior
  before the Bernoulli draw, so that the finite number of repetitions is
  carried through rather than assumed away. This is the headline at the 29-task
  suite scale (the $\sigma$ of the footnote in Section~\ref{sec:noise}).
\item \textbf{Plug-in distribution.} The same, with each measured rate treated
  as exact. The reweighted 117-task width and the same-system difference threshold ---
  $1.96\sqrt{2}\,\sigma$ of this distribution --- derive from it. Excluding
  parameter uncertainty makes both lower bounds on the re-run width, which is
  the conservative direction for reading a gap.
\item \textbf{Estimate bootstrap.} How precisely \emph{this experiment}
  determined its own number, by resampling repetitions within each task. It
  answers a question about our precision, not about a future run.
\end{itemize}

\section{Silent restarts}
\label{app:retries}
Across the 290 runs of the arm the harness's in-task retry fired 110 times.
Every one is instability that was absorbed before it could reach a score, which
is why the variance reported in Section~\ref{sec:noise} is a floor. The
concentration is the diagnostic: a hard task is hard in every round, whereas an
incident happens in one.

\begin{center}\small
\begin{tabular}{lcccccccccc}
\toprule
round & 1 & 2 & 3 & 4 & 5 & 6 & 7 & 8 & 9 & 10 \\
silent restarts & 6 & 5 & 4 & 3 & 3 & \textbf{68} & 3 & 7 & 7 & 4 \\
\bottomrule
\end{tabular}

\end{center}

\noindent By task, counting only tasks that were restarted at least once:

\begin{center}\scriptsize
\begin{tabular}{lc}
\toprule
task & silent restarts \\
\midrule
\texttt{MastodonImportMutedUsersTask} & 13 \\
\texttt{MastodonServerInfoReportTask} & 12 \\
\texttt{MastodonShareLocationTask} & 12 \\
\texttt{AcceptMeetingTask} & 9 \\
\texttt{CVEmailTask} & 9 \\
\texttt{ChangeWallpaperTask} & 9 \\
\texttt{MattermostCreateChannelTask} & 9 \\
\texttt{MattermostEmailTask} & 9 \\
\texttt{ScheduleLunchViaSmsTask} & 9 \\
\texttt{SendFormsTask} & 9 \\
\texttt{MattermostBudgetApprovalPipelineTask} & 5 \\
\texttt{MastodonGetServerInfoTask} & 4 \\
\texttt{DownloadSendReceiptTask} & 1 \\
\bottomrule
\end{tabular}

\end{center}

\section{Method and provenance}
\label{app:provenance}
The harness records a trajectory, a screenshot per step and a free-text
verdict, but no timestamps, model parameters, image digest or task-set commit.
Our orchestrator adds those, plus the count of directories the harness leaves
behind when it silently restarts a run, and writes one \texttt{fsync}'d line
per run. The raw record (\texttt{runs.jsonl}, 314 lines covering this arm and
the preparatory smoke and calibration arms) carries a sha256 that is stamped
into the analysis file and onto every report figure; the paper's figures
regenerate from that analysis file by one seeded command, so any figure can be
traced to the records that produced it.

Every field below is recorded on each of the 290 rows and checked for
consistency across them; a row disagreeing with the others would be reported as
a violation rather than averaged over.

\begin{center}\small
\begin{tabular}{ll}
\toprule
model & \texttt{moonshotai/kimi-k2.5} \\
serving provider / quantization & \texttt{atlas-cloud} / \texttt{int4} \\
sampling & temperature \texttt{0.0}, seed \texttt{42} \\
agent & \texttt{agents/pinned\_e2e\_agent.py} \\
image & \texttt{ghcr.io/tongyi-mai/mobile\_world:v1.4} \\
image digest & \begin{tabular}[t]{@{}l@{}}\texttt{sha256:00e24a8995892af9a0e1b68bc14ce455}\\\texttt{fae288d0f1ca8cf57e84a58e40e9f19e}\end{tabular} \\
task-set commit & \texttt{83e7b8fc75ebb6c4a098254a42999db5f4172666} \\
max rounds / step wait & 50 / 3.0s \\
harness auto-retry & 10 \\
container recycling & trial \\
host kernel / cores & 6.17.0-1022-gcp / 8 \\
pass threshold & score $>$ 0.99 \\
resampling & 10000 draws, seed \texttt{20260824} \\
rows analysed (headline basis) & 174 (rounds 4--10 less round 6; 203 over all seven, 290 over all ten) \\
duplicate trial keys & 0 (none expected) \\
provenance consistent across rows & yes \\
\bottomrule
\end{tabular}

\end{center}

\noindent Round boundaries in wall-clock time, so that the change of
Section~\ref{sec:endpoint} and the outage of Section~\ref{sec:noise} can
be located independently of the prose:

\begin{center}\small
\begin{tabular}{cccc}
\toprule
round & start (UTC) & end (UTC) & hours into the run \\
\midrule
1 & 08-25 20:10 & 08-25 23:17 & 0.0 \\
2 & 08-25 23:19 & 08-26 03:11 & 3.2 \\
3 & 08-26 03:13 & 08-26 06:47 & 7.1 \\
4 & 08-26 06:49 & 08-26 09:53 & 10.7\;$\leftarrow$ change \\
5 & 08-26 09:55 & 08-26 13:05 & 13.8 \\
6 & 08-26 13:07 & 08-26 16:25 & 17.0\;$\leftarrow$ outage \\
7 & 08-26 16:27 & 08-26 19:24 & 20.3 \\
8 & 08-26 19:26 & 08-26 22:18 & 23.3 \\
9 & 08-26 22:21 & 08-27 01:30 & 26.2 \\
10 & 08-27 01:33 & 08-27 04:51 & 29.4 \\
\bottomrule
\end{tabular}

\end{center}

\noindent The evaluation, the analysis and the figures are three commands:

\begin{quote}\ttfamily\small
python3 scripts/run\_matrix.py -{}-arm A\_temp0 -{}-trials 10 \\
\hspace*{2em}-{}-tasks data/pilot\_all.txt -{}-agent-type <pinned agent> \\
\hspace*{2em}-{}-temperature 0 -{}-seed 42 \\[2pt]
python3 scripts/analyze.py -{}-arm A\_temp0 -{}-trials 4:10 \\
\hspace*{2em}-{}-exclude-trials 6 -{}-leave-one-out -{}-regime-split 4 \\[2pt]
python3 scripts/plots\_paper.py
\end{quote}

\noindent Resampling is seeded, so the same raw records always produce the same
analysis file, and the analysis file carries the sha256 of the records it was
computed from. A separate script re-checks the paper's headline numeric claims against those
files and fails if any disagrees, or if the raw-record hash appears in the
text.

\noindent One class of number does not come from that pipeline: the endpoint
probes of Section~\ref{sec:determinism} and evidence 4--5 of
Section~\ref{sec:endpoint} were produced by a standalone diagnostic against the
live router and transcribed into a data file, which records that provenance
explicitly. We flag the difference rather than let it pass as equivalent.

\section{Tasks that randomise their own initial state}
\label{app:randomised}
Seven tasks in the benchmark call the standard library's random functions in
their own setup, and no seeding call appears anywhere in the codebase. For
these, successive runs are not literally the same problem, so a flip carries a
different meaning: the content changed, not only the outcome.

One of the seven, \texttt{CVEmailTask}, is in our task set. \textbf{It is
included in the 29-task aggregate}, and it flips (3/6 on the headline basis).
We flag it rather than drop it: excluding a task after seeing that it flips
would be selecting on the outcome, and one task cannot move a 29-task total by
more than 3.4 points in any case. The analysis records it as a separate
subgroup so that a reader who prefers to discount it can, and
Appendix~\ref{app:pertask} identifies it by name.

\section{Checking the alternative explanation: verifier errors}
\label{app:checker}
If the flips in Section~\ref{sec:noise} were scoring mistakes rather than
variance, the paper's central quantity would be measuring the wrong thing. We
therefore looked, and we had every incentive to find something: a demonstrable
verifier defect would have been a stronger and more quotable result than the
one we have, and it would have fitted this paper's theme better.

We ranked runs by how likely they were to be scoring errors --- minority
outcomes inside unstable tasks first, then passes completed in fewer steps than
any published run of the same task needed, then runs stopped at the round
budget. The top of that ranking led to one task; we added five whose failure
pattern looked wrong: an identical or implausible failure reason on every
trial, or a flip among tasks every published run passes. We inspected these six
tasks against the task definitions in the public repository.

Among the six we found no false positives and no false negatives. Four were
genuine agent failures (MattermostEmailTask, MastodonGetServerInfoTask,
MastodonShareLocationTask, CheckEventTimeTask). One was a task defect:
ReadQwen3PaperTask5 requires an exact match to \texttt{vie Latn,khm Khmr}, and
the English answer, a reversed order, and the paper's own spelling
\texttt{vie\_Latn} all score zero. Our ten runs all failed earlier, with empty
answers, so the defect did not affect our measurements; it is listed in
Table~\ref{tab:zero}. One remains unresolved: GraduationMassEmailTask, where
published data disagree, a published Kimi-K2.5 entry passes, and our ten runs
all failed with the same ``No email found'' error, which we did not pursue
further.

We report the result for three reasons: it is the obvious alternative
explanation and a reader deserves it answered before asking; it leaves that
explanation unsupported where it was most likely to show, which is what
Section~\ref{sec:noise} needs --- six tasks is a triage depth, not an
exhaustive audit, and the ranking is released so a reader can extend the search
from the top; and reporting a search that failed is the discipline this paper
argues for.
